\section{Un flux optionnel}
\label{sec:U5}

L'unité~4 a montré qu'un flux aléatoire se valorise en pondérant chaque issue par
sa probabilité. On ajoute ici le dernier ingrédient avant le contrat d'assurance
complet : le montant versé ne dépend plus seulement du hasard, mais d'un
\emph{choix} exercé au vu de ce hasard.

Ce déplacement est plus profond qu'il n'y paraît. Il invalide la méthode
précédente : espérer les issues sous leur probabilité réelle donnera désormais
un résultat faux. La raison de cet échec, et sa réparation, sont l'objet de
cette unité --- et le prérequis direct des modèles stochastiques de l'unité~8.

\subsection{La notion : payoff et asymétrie}

\begin{notion}{Un flux qui ne peut pas être négatif}
Une option d’achat (\emph{call}\index{call}) au prix d'exercice $K$ donne le droit d'acheter un \textit{sous-jacent} à ce prix $K$. Toutefois cet achat n’est pas obligatoire : le détenteur de l’option ne l’exerce que si cela lui profite. Le flux terminal, ou \emph{payoff}, vaut donc
\[
  C_T \;=\; \max(S_T - K,\; 0),
\]
et symétriquement pour une option de vente (\emph{put}\index{put}), droit de vendre, avec un flux terminal :
$P_T = \max(K - S_T,\; 0)$.

\medskip
L'opérateur $\max(\cdot,0)$ est tout ce qui distingue cette unité des
précédentes, et il change tout. La relation entre le sous-jacent et le flux
cesse d'être linéaire : les issues défavorables sont tronquées, les favorables
conservées. C'est l'\emph{asymétrie}\index{asymétrie}, et c'est très exactement
la forme d'une garantie d'assurance.
\end{notion}

\begin{center}
\begin{tikzpicture}[x=0.055cm, y=0.045cm]
  \draw[timeaxis] (0,0) -- (135,0);
  \draw[timeaxis] (0,-12) -- (0,58);
  \node[etiq, right] at (120,-9) {$S_T$};
  \node[etiq, above] at (0,58) {payoff};
  \draw[dashed, grisactu] (60,0) -- (60,-9);
  \node[etiq, below] at (60,-9) {$K$};
  \draw[very thick, vertactu] (0,0) -- (60,0) -- (120,52);
  \node[etiq, above left, vertactu] at (120,52) {call};
  \draw[very thick, rougeactu] (8,45) -- (60,0) -- (120,0);
  \node[etiq, above right, rougeactu] at (8,45) {put};
\end{tikzpicture}
\end{center}

\begin{exercice}[lecture d'un payoff]
Un call de prix d'exercice $K = 100$ est acheté $7{,}50$~\euro. Établir le payoff
puis le résultat net pour l'acheteur si le sous-jacent vaut à l'échéance
$88$, $100$, $104$ ou $121$. Déterminer le point mort.
\end{exercice}

\begin{corrige}
\begin{center}\small
\begin{tabular}{rrrr}
\toprule
$S_T$ & Payoff & Prime & Résultat net \\
\midrule
$88$  & $0$  & $-7{,}50$ & $-7{,}50$ \\
$100$ & $0$  & $-7{,}50$ & $-7{,}50$ \\
$104$ & $4$  & $-7{,}50$ & $-3{,}50$ \\
$121$ & $21$ & $-7{,}50$ & $+13{,}50$ \\
\bottomrule
\end{tabular}
\end{center}
Le point mort s'établit à $S_T = K + 7{,}50 = 107{,}50$.

\medskip
La perte est plafonnée à la prime, le gain ne l'est pas. Une compagnie qui
\emph{vend} une garantie occupe la position exactement inverse : gain plafonné à
la prime encaissée, perte non bornée. C'est la structure économique de tout
contrat d'assurance.
\end{corrige}

\begin{exercice}[pourquoi l'espérance réelle échoue]
Un actif vaut $100$ aujourd'hui. Dans un an, il vaudra $120$ avec une probabilité
de $60\,\%$, ou $90$ avec une probabilité de $40\,\%$. Le taux sans risque est de
$3\,\%$.

\textbf{(a)} Calculer l'espérance actualisée du sous-jacent selon la méthode de
l'unité~4. Comparer au prix observé, $100$. Que constate-t-on ?

\textbf{(b)} Un call de prix d'exercice $K=100$ est écrit sur cet actif. Calculer
son espérance actualisée par la même méthode.
\end{exercice}

\begin{corrige}
\textbf{(a)} $\mathbb{E}[S_1] = 0{,}6 \times 120 + 0{,}4 \times 90 = 108$, d'où
\[
  \frac{108}{1{,}03} = 104{,}85 \;\neq\; 100 .
\]
La méthode de l'unité~4 ne reproduit pas le prix de marché. L'écart n'est pas une erreur : il rémunère le risque supporté par le détenteur de l'actif. Un actif risqué ne s'actualise pas au taux sans risque.

\medskip
\textbf{(b)} Payoffs : $\max(120-100,0) = 20$ et $\max(90-100,0) = 0$.
\[
  \frac{0{,}6 \times 20 + 0{,}4 \times 0}{1{,}03} = \frac{12}{1{,}03} = 11{,}65 .
\]
Ce résultat est \emph{faux}, et pour la même raison qu'en (a) : il actualise un
flux risqué au taux sans risque. Mais il est ici impossible de le corriger par
une prime de risque ad hoc --- le risque d'une option n'est pas celui de son
sous-jacent, et il change à chaque instant. Il faut une autre méthode.
\end{corrige}

\subsection{Généralisation : réplication et repondération}

\begin{generalisation}{Ne pas prévoir, mais reproduire}
\textbf{L'idée.} Plutôt que de prédire la valeur de l'option, on construit un
portefeuille qui verse \emph{exactement} les mêmes flux dans tous les états de la nature.
Par absence d'opportunité d'arbitrage --- principe posé à l'unité~2 --- l'option
et ce portefeuille ont nécessairement le même prix.

\medskip
\textbf{Le modèle à une période.} Le sous-jacent vaut $S_0$ et prendra la valeur
$S_u = u\,S_0$ ou $S_d = d\,S_0$, avec $d < 1+r < u$. On cherche un portefeuille
composé de $\Delta$ unités de sous-jacent et de $B$ placé au taux sans risque,
qui reproduise les flux terminaux $(C_u, C_d)$. Ce portefeuille doit donc être tel que :\\
\begin{description}[leftmargin=3cm, style=nextline, font=\normalfont]
  \notation{\text{(cas de hausse)}}{$\Pi_u = \Delta S_u + B(1+r) = C_u$}
  \notation{\text{(cas de baisse)}}{$\Pi_d = \Delta S_d + B(1+r) = C_d$}
\end{description}

En résolvant, on trouve :\\

\[
  \Delta \;=\; \frac{C_u - C_d}{S_u - S_d},
  \qquad
  B \;=\; \frac{C_u - \Delta\,S_u}{1+r},
\]

L’hypothèse d’absence d'opportunité d'arbitrage – principe posé à l'unité 2 — conduit à la conclure que l'option et ce portefeuille ont nécessairement le même prix :\\
\[
  C_0 \;=\; \Delta\,S_0 + B .
\]


\medskip
\textbf{La repondération.} En réorganisant ces trois relations, le prix se
réécrit sous une forme remarquable :
\[
  C_0 \;=\; \frac{q\,C_u + (1-q)\,C_d}{1+r},
  \qquad\text{avec}\qquad
  q \;=\; \frac{(1+r) - d}{u - d}.
\]
Le prix est donc une espérance actualisée, comme à l'unité~4 --- mais calculée
avec $q$ au lieu de la probabilité réelle $p$.

\medskip
\textbf{Ce qu'est $q$, et ce qu'il n'est pas.} La quantité $q$, dite
\emph{probabilité risque-neutre}\index{probabilité risque-neutre}, est comprise
entre $0$ et $1$ et se manipule comme une probabilité. Mais elle ne décrit
\emph{aucune} croyance sur l'avenir : elle ne dépend pas de $p$, et $p$
n'apparaît nulle part dans le prix. C'est un artefact de calcul, la pondération
qui rend l'actualisation au taux sans risque légitime.

\medskip
\textbf{L'extension.} Sur plusieurs périodes, on répète l'opération de proche en
proche, en partant des dates terminales et en remontant. Lorsque le nombre de
périodes devient très grand, cette récurrence converge vers une intégrale : c'est
la formule de Black-Scholes. Lorsque les branches deviennent trop nombreuses pour
être énumérées, on tire au sort des trajectoires sous $q$ et l'on moyenne : c'est
la \emph{simulation de Monte-Carlo}\index{Monte-Carlo}.
\end{generalisation}

\begin{center}
\begin{tikzpicture}[x=1.6cm, y=0.95cm]
  \node[etiq] (a) at (0,0) {$S_0 = 100$};
  \node[etiq] (u) at (1.6,1.15) {$120$};
  \node[etiq] (d) at (1.6,-1.15) {$90$};
  \draw[flux] (a) -- (u) node[midway, above left, etiq] {$q$};
  \draw[fluxneg] (a) -- (d) node[midway, below left, etiq] {$1-q$};
  \node[etiq, right] at (2.0,1.15) {$C_u = 20$};
  \node[etiq, right] at (2.0,-1.15) {$C_d = 0$};
  \draw[dashed, grisactu] (0,-2.0) -- (0,-2.4) node[below, etiq] {$t=0$};
  \draw[dashed, grisactu] (1.6,-2.0) -- (1.6,-2.4) node[below, etiq] {$t=1$};
\end{tikzpicture}
\end{center}

\begin{exercice}[réplication et repondération]
Reprendre les données de l'exercice~5.2 : $S_0 = 100$, $u = 1{,}20$,
$d = 0{,}90$, $r = 3\,\%$, call de prix d'exercice $K = 100$.

\textbf{(a)} Construire le portefeuille de réplication et en déduire le prix.

\textbf{(b)} Calculer $q$ et retrouver le prix par la formule de repondération.

\textbf{(c)} Que devient le prix si la probabilité réelle de hausse passe de
$60\,\%$ à $80\,\%$ ?
\end{exercice}

\begin{corrige}
\textbf{(a)} Payoffs : $C_u = 20$, $C_d = 0$.
\[
  \Delta = \frac{20-0}{120-90} = 0{,}6667,
  \qquad
  B = \frac{20 - 0{,}6667 \times 120}{1{,}03} = \frac{-60}{1{,}03} = -58{,}25 .
\]
\[
  C_0 = 0{,}6667 \times 100 - 58{,}25 = 8{,}41 .
\]
Le portefeuille consiste à détenir deux tiers d'unité de sous-jacent, financés
pour $58{,}25$ par emprunt. Vérification dans l'issue basse :
$0{,}6667 \times 90 - 58{,}25 \times 1{,}03 = 60 - 60 = 0 = C_d$.

\medskip
\textbf{(b)} $q = \dfrac{1{,}03 - 0{,}90}{1{,}20-0{,}90} = \dfrac{0{,}13}{0{,}30} = 0{,}4333$.
\[
  C_0 = \frac{0{,}4333 \times 20 + 0{,}5667 \times 0}{1{,}03} = \frac{8{,}6667}{1{,}03} = 8{,}41 .
\]
Les deux méthodes coïncident, comme elles le doivent.

\medskip
\textbf{(c)} Le prix ne change pas : il reste $8{,}41$. La probabilité réelle
n'entre dans aucune des deux formules.

\medskip
Ce résultat déconcerte, et c'est normal. Il signifie qu'anticiper une hausse ne
rend pas l'option plus chère : ce qui la valorise n'est pas l'espérance de gain,
mais le coût de sa reproduction. Noter au passage l'écart avec le calcul erroné
de l'exercice~5.2 : $8{,}41$ contre $11{,}65$, soit une surestimation de $38{,}5\,\%$.
\end{corrige}

\begin{exercice}[garantie plancher sur un contrat d'épargne]
Un contrat d'épargne investit $10\,000$~\euro{} sur un support dont la valeur
peut, chaque année, progresser de $15\,\%$ ou reculer de $10\,\%$. Le taux sans
risque est de $2\,\%$. Le contrat garantit à l'assuré de récupérer au minimum sa
mise au terme de deux ans.

\textbf{(a)} Construire l'arbre des valeurs du support sur deux périodes.

\textbf{(b)} Valoriser la garantie, c'est-à-dire l'engagement de l'assureur de
combler l'écart si la valeur finale est inférieure à $10\,000$~\euro.

\textbf{(c)} Commenter ce que représente ce montant pour l'assureur.
\end{exercice}

\begin{corrige}
\textbf{(a)} Avec $u = 1{,}15$ et $d = 0{,}90$ :
\begin{center}\small
\begin{tabular}{lrr}
\toprule
Trajectoire & Valeur en $t=2$ & Coût de la garantie \\
\midrule
$uu$ & $13\,225$ & $0$ \\
$ud$ et $du$ & $10\,350$ & $0$ \\
$dd$ & $8\,100$ & $1\,900$ \\
\bottomrule
\end{tabular}
\end{center}

\medskip
\textbf{(b)} $q = \dfrac{1{,}02-0{,}90}{1{,}15-0{,}90} = \dfrac{0{,}12}{0{,}25} = 0{,}48$.

La garantie ne se déclenche que sur la trajectoire $dd$, de probabilité
risque-neutre $(1-q)^2 = 0{,}52^2 = 0{,}2704$ :
\[
  V_0 = \frac{0{,}2704 \times 1\,900}{(1{,}02)^{2}}
      = \frac{513{,}76}{1{,}0404} \approx 493{,}81\ \text{\euro}.
\]

\medskip
\textbf{(c)} Ce montant, soit $4{,}94\,\%$ de la mise, est ce que coûterait à
l'assureur la couverture de sa promesse sur le marché. Il ne s'agit pas d'une
perte probable, mais du prix d'un engagement.

C'est très exactement la nature de la \emph{valeur temps des options et
garanties} incluse dans une meilleure estimation : un contrat d'épargne à
garantie plancher vaut davantage que la seule projection de son support, et
l'écart est une option que l'assureur a vendue sans toujours l'avoir facturée.
\end{corrige}

\begin{pont}{Où cela sert dans Solvabilité~II}
Les contrats d'assurance sont truffés d'options, le plus souvent implicites et
détenues par l'assuré : garantie de taux minimum, participation aux bénéfices,
faculté de rachat, option de conversion en rente. Toutes ont la forme
asymétrique de cette unité, et toutes sont exercées au moment le moins favorable
à l'assureur.

\medskip
Trois conséquences directes.

\emph{La meilleure estimation ne peut pas ignorer ces options.} La directive
impose de tenir compte de la valeur des garanties financières et des options
contractuelles. Une projection en scénario central les manquerait par
construction : dans un scénario moyen, une garantie plancher ne se déclenche
jamais.

\emph{C'est pourquoi les modèles vie sont stochastiques.} Là où un portefeuille
non-vie se projette sur une cadence déterministe, un portefeuille épargne exige
de simuler des milliers de trajectoires --- sous la probabilité $q$ de cette
unité --- et de moyenner. L'unité~8 en fait la mise en \oe uvre.

\emph{Les hypothèses comportementales redeviennent réelles.} Le taux de rachat
des assurés, lui, n'est pas une probabilité risque-neutre : c'est une hypothèse
d'expérience au sens de l'unité~4. Un modèle vie fait donc coexister les deux
natures de probabilité, et les confondre est l'une des erreurs les plus
fréquentes.
\end{pont}

\begin{refreg}
\textbf{Références.} Directive 2009/138/CE, article~77, paragraphe~2 : le calcul
de la meilleure estimation tient compte de la valeur des garanties financières et
de toute option contractuelle incluse dans les contrats. \\
Article~79 : évaluation des garanties financières et des options contractuelles,
et exigence que les hypothèses relatives à leur exercice soient réalistes et
fondées sur des informations actuelles et crédibles. \\
Règlement délégué (UE) 2015/35 : méthodes de valorisation et cohérence avec les
données de marché observables. \emph{[Références d'articles à préciser avant
diffusion.]}
\end{refreg}

\subsection*{Notations introduites}
\begin{description}[leftmargin=2.4cm, style=nextline, font=\normalfont]
  \notation{K}{prix d'exercice}
  \notation{S_T}{valeur du sous-jacent à l'échéance}
  \notation{u,\ d}{facteurs de hausse et de baisse}
  \notation{\Delta}{quantité de sous-jacent du portefeuille de réplication}
  \notation{B}{montant placé au taux sans risque}
  \notation{q}{probabilité risque-neutre}
\end{description}
